problem:  
Let \((X,d,\mu)\) be an \(\mathrm{RCD}(K,\infty)\) metric–measure space with \(K>0\). Denote by \(\lambda(X)\) the optimal logarithmic Sobolev constant, i.e. the largest number \(\lambda>0\) such that for all \(u\in W^{1,2}(X)\) with \(\int u^2\,d\mu =1\),
\[
\int_X u^2\log u^2\,d\mu \le \frac{2}{\lambda}\int_X |\nabla u|^2\,d\mu .
\]
It is known that \(\lambda(X)\ge 2K\) by the Bakry–Émery calculus extended to the \(\mathrm{RCD}(K,\infty)\) setting.  

**Problem:**  
Develop a rigidity–stability dichotomy for the logarithmic Sobolev constant in the nonsmooth \(\mathrm{RCD}(K,\infty)\) context:

1. (**Rigidity**) If \(\lambda(X)=2K\), must \((X,d,\mu)\) be isomorphic (in the measured Gromov sense) to a Gaussian Euclidean space \((\mathbb{R}^n,|\cdot|,\gamma_K)\) or to a metric product of such a Gaussian space with an \(\mathrm{RCD}(K,\infty)\) space of zero entropy growth?  
2. (**Quantitative stability**) If \(\lambda(X)\) is close to \(2K\), can the deviation \(\lambda(X)-2K\) be bounded below by a quantitative measure of “non-Gaussianness,” such as the Wasserstein distance between \(\mu\) and a corresponding Gaussian measure, or by a curvature deficit functional derived from the Γ₂-calculus?

Why is it a "good" problem:  
This problem pushes the classical Bakry–Émery theory into the genuinely nonsmooth regime of \(\mathrm{RCD}(K,\infty)\) spaces, where the analytic and geometric structures interact through optimal transport rather than differential tensors. Determining whether Gaussian rigidity and quantitative stability persist in this setting would unify geometric measure theory, optimal transport, and functional inequalities. It demands new methods to capture curvature and entropy in non-smooth contexts and aims to establish a sharp rigidity/stability threshold for the fundamental log-Sobolev inequality. The statement is simple—asking whether equality forces Gaussian geometry—but its depth is immense: success would extend one of the deepest links between curvature and functional inequalities beyond the smooth manifold framework.